\documentclass[11pt]{article}
\usepackage[T1]{fontenc}
\usepackage{lmodern,amsmath,amssymb,amsthm,mathtools,microtype}
\usepackage[margin=1.15in]{geometry}
\usepackage[hidelinks]{hyperref}
\newtheorem{theorem}{Theorem}[section]
\newtheorem{proposition}[theorem]{Proposition}
\newtheorem{lemma}[theorem]{Lemma}
\theoremstyle{remark}
\newtheorem{remark}[theorem]{Remark}
\title{The Level-Six Fibre Lattice in Engel's Torus Fibration}
\author{Raeez Lorgat}
\date{}
\hypersetup{
  pdftitle={The Level-Six Fibre Lattice in Engel's Torus Fibration},
  pdfauthor={Raeez Lorgat},
  pdfsubject={Integral monodromy, discriminant pairings, and two-fibre topology}
}
\begin{document}
\maketitle

\begin{abstract}
For the lattice $\Lambda=H_1(F,\mathbb Z)$ of a smooth complex two-torus
fibre in Engel's fibration, we recover the two integral monodromy markings
and determine the unique primitive invariant alternating form $\xi$.  The
distinguished circle class $\delta_0$ and invariant covector $\psi$ satisfy
\[
                 \iota_{\delta_0}\xi=-6\psi.
\]
An explicit unimodular basis puts $\xi$ in the form $J\oplus6J$; hence its
discriminant group is $(\mathbb Z/6\mathbb Z)^2$, with a perfect alternating
pairing and a six-dimensional irreducible finite Heisenberg representation.
We also compute the complete Smith form of a two-fibre presentation,
including the zero-determinant case, and identify its determinant with a
denominator-cleared Seifert Euler number.  The Leray transgression is stated
only for Engel's multiplicities $(3,4)$; no transgression formula for
arbitrary multiplicities is asserted.
\end{abstract}

\section{The fibre calculation}

The calculation to be explained is the contraction identity
\[
                         \iota_{\delta_0}\xi=-6\psi.       \tag{1.1}
\]
Its content is not the coefficient alone: $\delta_0$, $\xi$, and $\psi$
must be carried by the same integral geometric lattice, with compatible
markings and orientations.

Let $f:Y(u)\to\mathbb P^1$ be the complex two-torus fibration described by
Engel, and put
\[
 U=\mathbb P^1\setminus\{0,1,\infty\},\qquad
 F=f^{-1}(t),\qquad \Lambda=H_1(F,\mathbb Z),\qquad
 V=\Lambda^\vee
\]
for $t\in U$.  Over $U$ this is also the smooth fibre lattice of the
unmodified fibration $X(u)\to\mathbb P^1$.  Engel's construction supplies
the filtration
\[
 0\longrightarrow\Lambda^\times\longrightarrow\Lambda
   \mathrel{\mathop{\longrightarrow}^{\psi}}\mathbb Z\longrightarrow0,
 \qquad
 0\longrightarrow\mathbb Z\delta_0\longrightarrow\Lambda^\times
   \longrightarrow\Lambda'\longrightarrow0,                 \tag{1.2}
\]
where $\delta_0$ is the circle class of the principal
$\mathbb G_m$-bundle in Engel's chosen marking
\cite[Sections 2 and 4]{Engel}.

We record the integral markings because every sign below depends on them.
Let $\varrho=e^{2\pi i/3}$.  In the two good-reduction markings, and writing
$\tau_0,\tau_1$ for the periods of the second elliptic factors, Engel takes
\[
\begin{aligned}
 e_1&=(1,0),& e_2&=(\varrho,0),& \delta_0&=(0,1),&e_4&=(0,\tau_0),
 &e_3&=\tfrac13(2+\varrho,1),\\
 f_1&=(1,0),& f_2&=(i,0),& \delta_1&=(0,1),&f_4&=(0,\tau_1),
 &f_3&=\tfrac12(1+i,1).
\end{aligned}                                                   \tag{1.3}
\]
The first four entries in each row generate a direct-sum sublattice; after
adjoining $e_3$, respectively $f_3$, the resulting overlattice $\Lambda$ has
ordered bases
\[
             B_0=(e_1,e_2,e_3,e_4),\qquad
             B_1=(f_1,f_2,f_3,f_4).
\]
For standard counterclockwise meridians
$\gamma_0\gamma_1\gamma_\infty=1$, the matrices act on column coordinates
and are
\[
 [T_0]_{B_0}=
 \begin{pmatrix}
 0&-1&-1&0\\ 1&-1&0&0\\ 0&0&1&0\\ 0&0&0&1
 \end{pmatrix},\qquad
 [T_1]_{B_1}=
 \begin{pmatrix}
 0&1&0&0\\ -1&0&-1&0\\ 0&0&1&0\\ 0&0&0&1
 \end{pmatrix}.                                                  \tag{1.4}
\]
The complex orientations give multiplication by $\varrho$ at $0$ and by
$-i$ at $1$.  The circle classes and the transition matrix are
\[
 \delta_0=-2e_1-e_2+3e_3,\qquad
 \delta_1=-f_1-f_2+2f_3,\qquad
 C=
 \begin{pmatrix}
 1&1&0&1\\0&1&0&0\\0&-1&1&0\\0&0&0&1
 \end{pmatrix},                                                   \tag{1.5}
\]
where the columns of $C$ express $B_1$ in $B_0$ and
$C\delta_1=\delta_0$.  These are precisely the integral data of
Engel's Theorem~4.1.  Its proof fixes the sign at infinity by taking a
simple zero, rather than a pole, of the linearization for a counterclockwise
meridian \cite[pp.~10--13]{Engel}.

In the common basis $B_0$, direct conjugation and
$T_0T_1T_\infty=I$ give
\[
 [T_1]_{B_0}=C[T_1]_{B_1}C^{-1}=
 \begin{pmatrix}
 -1&1&-1&2\\-1&0&-1&1\\1&1&2&-1\\0&0&0&1
 \end{pmatrix},                                                   \tag{1.6}
\]
\[
 [T_\infty]_{B_0}=
 \begin{pmatrix}
 2&1&1&0\\0&1&0&-1\\-1&-1&0&1\\0&0&0&1
 \end{pmatrix}.                                                   \tag{1.7}
\]
Thus $T_0^3=T_1^4=I$, $(T_\infty-I)^2=0$, and
\[
 (T_\infty-I)\Lambda
   =\langle e_1-e_3,e_3-e_2\rangle.                               \tag{1.8}
\]

\begin{theorem}[The geometric level-six lattice]
Let $\Gamma=\langle T_0,T_1\rangle$.  For the fibre lattice and markings
above,
\[
 \Lambda^\Gamma=\mathbb Z\delta_0,
 \qquad V^\Gamma=\mathbb Z\psi,
 \qquad (\textstyle\bigwedge^2 V)^\Gamma=\mathbb Z\xi,             \tag{1.9}
\]
where $\psi=e_4^\vee$ and
\[
 \xi=3e_1^\vee\wedge e_2^\vee+e_1^\vee\wedge e_3^\vee
      -2e_2^\vee\wedge e_3^\vee-2e_3^\vee\wedge e_4^\vee.       \tag{1.10}
\]
With the contraction convention
$(\iota_v\xi)(w)=\xi(v,w)$, identity \emph{(1.1)} holds.  Moreover, the
ordered basis
\[
 a_1=e_1,\qquad b_1=e_3,\qquad
 a_2=-\delta_0=2e_1+e_2-3e_3,\qquad b_2=e_4-2e_1             \tag{1.11}
\]
is unimodular and satisfies
\[
 [\xi]_{(a_1,b_1,a_2,b_2)}=J\oplus6J,\qquad
 J=\begin{pmatrix}0&1\\-1&0\end{pmatrix},\qquad
 \delta_0=-a_2,\quad\psi=b_2^\vee.                             \tag{1.12}
\]
Consequently
\[
 \det(\xi^\flat)=36,\qquad
 \operatorname{SNF}(\xi^\flat)=\operatorname{diag}(1,1,6,6),     \tag{1.13}
\]
and the elementary divisors of $\xi$ are $(1,6)$.
\end{theorem}

\begin{proof}
Solving $T_iv=v$ for $i=0,1$ gives
$v=q(-2,-1,3,0)^T$, and solving $T_i^{-T}u=u$ gives
$u=q(0,0,0,1)^T$.  Integrality makes the two common invariant lattices
exactly $\mathbb Z\delta_0$ and $\mathbb Z\psi$.

Write the coefficients of an alternating form in the order
\[
 e_1^\vee\wedge e_2^\vee, e_1^\vee\wedge e_3^\vee,
 e_1^\vee\wedge e_4^\vee, e_2^\vee\wedge e_3^\vee,
 e_2^\vee\wedge e_4^\vee, e_3^\vee\wedge e_4^\vee.
\]
Row reduction of $T_i^TMT_i=M$ gives the integral solutions
$q(3,1,0,-2,0,-2)$, proving the last equality in (1.9) and recovering
Engel's invariant class from Lemma~6.2 \cite[p.~15]{Engel}.  Its matrix in
$B_0$ is
\[
 [\xi]_{B_0}=
 \begin{pmatrix}
 0&3&1&0\\-3&0&-2&0\\-1&2&0&-2\\0&0&2&0
 \end{pmatrix}.                                                   \tag{1.14}
\]
Therefore
\[
       \delta_0^T[\xi]_{B_0}=(0,0,0,-6),
\]
which is (1.1).

The columns of (1.11) form the matrix
\[
 P=\begin{pmatrix}1&0&2&-2\\0&0&1&0\\0&1&-3&0\\0&0&0&1\end{pmatrix},
 \qquad \det P=-1.
\]
Multiplication gives $P^T[\xi]_{B_0}P=J\oplus6J$.  This proves the
normal form and the Smith form.  Finally,
\[
 \operatorname{Pf}([\xi]_{B_0})=-6,\qquad
 \tfrac12\xi\wedge\xi
   =-6e_1^\vee\wedge e_2^\vee\wedge e_3^\vee\wedge e_4^\vee,
\]
so $\det[\xi]_{B_0}=(-6)^2=36$, as asserted.
\end{proof}

\begin{remark}[Basis and orientation]
If the columns of $A\in\operatorname{GL}_4(\mathbb Z)$ are a new basis in
an old one, then the simultaneous coordinate changes are
\[
 [\xi]'=A^T[\xi]A,\qquad [\delta]'=A^{-1}[\delta],\qquad
 [\psi]'=A^T[\psi].                                               \tag{1.15}
\]
Thus (1.1) is tensorial.  The Pfaffian changes by $\det A$, which explains
the signs $-6$ in $B_0$ and $+6$ in the orientation-reversing basis
(1.11), while the determinant and Smith form do not change.  Reversing only
$\delta_0$, only $\psi$, or the contraction convention reverses the sign in
(1.1).  The abstract normal form alone does not select these classes: for
example, in $J\oplus6J$, taking $\delta=a_1$ and $\psi=b_2^\vee$ gives
$\iota_\delta\xi=b_1^\vee$, not $-6\psi$.
\end{remark}

\section{The discriminant pairing and Heisenberg module}

Let $(L,\xi)$ be any rank-four alternating lattice with the normal form
$J\oplus6J$ in a basis $(a_1,b_1,a_2,b_2)$.  Define
\[
 \xi^\flat:L\longrightarrow L^\vee,\qquad \xi^\flat(x)=\xi(x,-),
 \qquad
 \mathcal D_\xi=L^\vee/\xi^\flat(L).
\]
The rational extension
$\phi=\xi^\flat\otimes\mathbb Q$ is an isomorphism.

\begin{lemma}[Discriminant pairing]
For $u,v\in L^\vee$, the formula
\[
 b_\xi(\bar u,\bar v)=u(\phi^{-1}v)\pmod{\mathbb Z}               \tag{2.1}
\]
defines a perfect alternating pairing
$b_\xi:\mathcal D_\xi\times\mathcal D_\xi\to\mathbb Q/\mathbb Z$.
Moreover, if
\[
 \alpha=[a_2^\vee],\qquad \beta=[b_2^\vee],
\]
then
\[
 \mathcal D_\xi=\langle\alpha\rangle\oplus\langle\beta\rangle
 \cong(\mathbb Z/6\mathbb Z)^2,\qquad
 b_\xi(\alpha,\beta)=\tfrac16,\quad
 b_\xi(\beta,\alpha)=-\tfrac16.                                  \tag{2.2}
\]
\end{lemma}

\begin{proof}
Write $x=\phi^{-1}u$ and $y=\phi^{-1}v$.  Replacing $u$ by
$u+\xi^\flat(\lambda)$ changes (2.1) by
$\xi(\lambda,y)=-v(\lambda)\in\mathbb Z$; replacing $v$ by
$v+\xi^\flat(\mu)$ changes it by $u(\mu)\in\mathbb Z$.  The pairing is
therefore well defined and alternating.  If it annihilates $\bar u$ against
every $\bar v$, then $v(x)=-u(\phi^{-1}v)$ is integral for every
$v\in L^\vee$.  Hence $x\in L$, so $\bar u=0$.  The induced homomorphism
to the Pontryagin dual is injective between finite groups of the same order,
and is therefore an isomorphism.  Formula (2.2) follows directly from the
normal form.
\end{proof}

Put $K=\langle\alpha\rangle\cong\mathbb Z/6\mathbb Z$, let
$\zeta_6=e^{2\pi i/6}$, and let $\chi\in\widehat K$ satisfy
$\chi(\alpha)=\zeta_6$.  The map
\[
          r\alpha+s\beta\longmapsto(r\alpha,\chi^s)
\]
is a polarization $\mathcal D_\xi\cong K\oplus\widehat K$.  Define
\[
 H_6=\mu_6\times K\times\widehat K
\]
with multiplication
\[
 (z,u,\eta)(z',u',\eta')
   =\bigl(zz'\eta'(u),u+u',\eta\eta'\bigr).                       \tag{2.3}
\]

\begin{theorem}[Level-six Heisenberg module]
The center of $H_6$ is the displayed copy of $\mu_6$.  Its commutator
pairing is
\[
 [(1,u,\eta),(1,u',\eta')]
  =\bigl(\eta'(u)\eta(u')^{-1},0,1\bigr)
  =\bigl(e^{2\pi i b_\xi(d,d')},0,1\bigr),                        \tag{2.4}
\]
where $d\leftrightarrow(u,\eta)$ and $d'\leftrightarrow(u',\eta')$.
On the six-dimensional space $\mathcal H=\mathbb C[K]$, the formula
\[
 [\rho(z,u,\eta)h](t)=z\eta(t)h(t+u)                              \tag{2.5}
\]
defines an irreducible representation, and the central $\mu_6$ acts by its
faithful tautological character.
\end{theorem}

\begin{proof}
The bicharacter in (2.3) makes the multiplication associative, and direct
multiplication gives (2.4).  Its nondegeneracy forces a central element to
have $u=0$ and $\eta=1$.  The same calculation verifies (2.5).

For $t\in K$, let $\delta_t$ be the delta-function at $t$.  The operator
\[
 P_t=\frac16\sum_{\eta\in\widehat K}
        \eta(t)^{-1}\rho(1,0,\eta)
\]
projects onto $\mathbb C\delta_t$.  A nonzero invariant subspace therefore
contains some $\delta_t$, and translations carry it through every
delta-function.  These span $\mathcal H$, proving irreducibility.
\end{proof}

Applied to Theorem~1.1, this construction belongs to the fibre lattice
$H_1(F,\mathbb Z)$.  It is not constructed from the ordinary cohomology of
the total six-manifold.

\section{Two multiple fibres}

The topology of two logarithmic fillings leads to an elementary presentation
whose full integral content is visible before any Leray differential is
considered.  It is useful to include a global section integer from the start.

\begin{theorem}[Two-fibre Smith form]
Let $r,s>0$, let $m,n,b_0\in\mathbb Z$, and define
\[
 G_{b_0}=\left\langle c,x_0,x_1\ \middle|\
 c\text{ central},\ x_0x_1=c^{b_0},\ x_0^r=c^m,\ x_1^s=c^n
 \right\rangle.                                                   \tag{3.1}
\]
Put
\[
 \Delta_{b_0}=rn+sm-rsb_0,\qquad d=\gcd(r,s,m,n).                \tag{3.2}
\]
Then
\[
 G_{b_0}\cong\operatorname{coker}
 \begin{pmatrix}r&s\\-m&n-sb_0\end{pmatrix}.                     \tag{3.3}
\]
If $\Delta_{b_0}\ne0$, then
\[
 G_{b_0}\cong\mathbb Z/d\mathbb Z
       \oplus\mathbb Z/(|\Delta_{b_0}|/d)\mathbb Z,\qquad
 |G_{b_0}|=|\Delta_{b_0}|.                                        \tag{3.4}
\]
If $\Delta_{b_0}=0$, then
\[
                    G_{b_0}\cong\mathbb Z\oplus\mathbb Z/d\mathbb Z.
                                                                    \tag{3.5}
\]
In either case $G_{b_0}$ is cyclic if and only if $d=1$.
\end{theorem}

\begin{proof}
The product relation gives $x_1=x_0^{-1}c^{b_0}$, so the group is abelian.
In additive notation its two remaining relation vectors in
$\mathbb Zx_0\oplus\mathbb Zc$ are
\[
          rx_0-mc,\qquad sx_0+(n-sb_0)c,
\]
which proves (3.3).  The greatest common divisor of the entries is
\[
 \gcd(r,s,m,n-sb_0)=\gcd(r,s,m,n)=d,
\]
and the determinant is $\Delta_{b_0}$.  For nonzero determinant, the Smith
invariants are $d$ and $|\Delta_{b_0}|/d$.  If the determinant vanishes, the
matrix still has rank one because $r,s>0$, so its Smith form is
$\operatorname{diag}(d,0)$.  The cyclicity assertions follow.
\end{proof}

\begin{proposition}[Section gauges and orientations]
For $k_0,k_1\in\mathbb Z$, the generator change
\[
 x_0'=x_0c^{k_0},\qquad x_1'=x_1c^{k_1}                            \tag{3.6}
\]
sends
\[
 m'=m+rk_0,\qquad n'=n+sk_1,\qquad b_0'=b_0+k_0+k_1             \tag{3.7}
\]
and preserves both $d$ and $\Delta_{b_0}$.  In the gauge $b_0=0$, the
presentation-preserving choice $k_1=-k_0=-k$ sends
\[
                         (m,n)\longmapsto(m+rk,n-sk).              \tag{3.8}
\]
Interchanging $(r,m,x_0)$ with $(s,n,x_1)$ preserves
$\Delta_{b_0}$.  Replacing the oriented fibre generator by $c'=c^{-1}$
sends $(m,n,b_0)$ to $(-m,-n,-b_0)$ and sends
$\Delta_{b_0}$ to $-\Delta_{b_0}$.
\end{proposition}

\begin{proof}
Substitution in the three defining relations gives (3.7), and
\[
 r(n+sk_1)+s(m+rk_0)-rs(b_0+k_0+k_1)=\Delta_{b_0}.
\]
The remaining statements follow by the indicated changes of generators.
They also show why only $|\Delta_{b_0}|$ is intrinsic without fixed
orientations.
\end{proof}

The hypotheses in Theorem~3.1 cannot be weakened silently.  For $b_0=0$,
\[
 (r,s;m,n)=(2,4;2,2)\quad\Longrightarrow\quad
 G_0\cong\mathbb Z/2\mathbb Z\oplus\mathbb Z/6\mathbb Z,
\]
although $\Delta_0=12$.  Also
\[
 (2,4;2,-4)\quad\Longrightarrow\quad
 G_0\cong\mathbb Z\oplus\mathbb Z/2\mathbb Z,\qquad \Delta_0=0.
\]
At the simplest zero-coefficient boundary, $m=n=0$ gives
$G_0\cong\mathbb Z\oplus\mathbb Z/\gcd(r,s)\mathbb Z$.

\begin{proposition}[Seifert Euler number]
Let $M$ be an oriented Seifert manifold over $S^2$ with two normalized
effective exceptional fibres, written in Wilkes's convention as
\[
       (b_0,S^2;(r,\beta_0),(s,\beta_1)),                          \tag{3.9}
\]
where $r,s\ge2$, $0<\beta_0<r$, $0<\beta_1<s$, and
$\gcd(r,\beta_0)=\gcd(s,\beta_1)=1$.  With
\[
                 m=-\beta_0,\qquad n=-\beta_1,                   \tag{3.10}
\]
one has $\pi_1(M)\cong G_{b_0}$ and
\[
 e(M)=-\left(b_0+\frac{\beta_0}{r}+\frac{\beta_1}{s}\right)
      =-b_0+\frac mr+\frac ns,\qquad
 rs\,e(M)=\Delta_{b_0}.                                           \tag{3.11}
\]
In particular $d=1$, and hence
\[
 \pi_1(M)\cong
 \begin{cases}
  \mathbb Z/|\Delta_{b_0}|\mathbb Z,&e(M)\ne0,\\
  \mathbb Z,&e(M)=0.
 \end{cases}                                                       \tag{3.12}
\]
If $g=\gcd(r,s)$ and $L=\operatorname{lcm}(r,s)=rs/g$, then the smaller
denominator clearing is
\[
             L e(M)=-Lb_0+\frac{s}{g}m+\frac{r}{g}n
                    =\frac{\Delta_{b_0}}{g}\in\mathbb Z.          \tag{3.13}
\]
\end{proposition}

\begin{proof}
Wilkes's presentation has central regular fibre $h$, local relations
$a_i^{\alpha_i}h^{\beta_i}=1$, and, over $S^2$ with two exceptional
fibres, the product relation $a_0a_1=h^{b_0}$.  Setting
$c=h$, $x_i=a_i$ gives (3.1) and (3.10).  His Euler-number convention gives
(3.11) \cite[Proposition 4.3, p.~149]{Wilkes}.  Effectiveness implies
$d=1$, and Theorem~3.1 gives (3.12).  Equation (3.13) is immediate.
Hatcher's section-change calculation gives the geometric counterpart of
(3.6): twists are added at one boundary torus and subtracted at another,
leaving the Euler sum fixed \cite[Section 2.1, pp.~25--27]{Hatcher}.
\end{proof}

\begin{remark}
The rational number $e(M)$ is orientation-sensitive.  Reversing the fibre
orientation and renormalizing Wilkes's two local pairs sends
\[
 (b_0;\beta_0,\beta_1)
   \longmapsto(-b_0-2;r-\beta_0,s-\beta_1),
\]
which negates both $e(M)$ and $\Delta_{b_0}$ but preserves the abstract
group.  With exceptional fibres, $e(M)$ is a rational Seifert Euler number;
only in the ordinary circle-bundle case is it itself the integral Euler
class on the underlying base.
\end{remark}

\section{The Leray carrier in Engel's (3,4) case}

Choose lifts of Engel's logarithmic-transform classes
\[
 a_0\in\tfrac13\Lambda^{T_0},\qquad
 a_1\in\tfrac14\Lambda^{T_1},
\]
and put
\[
 v_0=3a_0,\qquad v_1=4a_1,\qquad
 m=\psi(v_0),\qquad n=\psi(v_1),\qquad
 p=12\psi(a_0+a_1)=4m+3n.                                        \tag{4.1}
\]
When the classes of $a_0,a_1$ are primitive torsion classes, the fibres at
$0$ and $1$ have multiplicities $3$ and $4$, respectively
\cite[Corollary 2.10]{Engel}.

The carrier of the transgression is the cohomological Leray spectral
sequence of the specific map $f:Y(u)\to\mathbb P^1$,
\[
 E_2^{a,b}=H^a(\mathbb P^1,R^bf_*\mathbb Z)
       \Longrightarrow H^{a+b}(Y(u),\mathbb Z).                   \tag{4.2}
\]
Over $U$, the restriction of $R^1f_*\mathbb Z$ is a local system with
fibre $V=H^1(F,\mathbb Z)$ and monodromy $T_i^{-T}$.  Across the three
singular values, $R^1f_*\mathbb Z$ is the corresponding constructible
sheaf, not a local system.  The fibres of $f$ are connected, so
$R^0f_*\mathbb Z\cong\mathbb Z_{\mathbb P^1}$.  Thus the first relevant
differential is precisely
\[
 d_2^{0,1}:H^0(\mathbb P^1,R^1f_*\mathbb Z)
       \longrightarrow H^2(\mathbb P^1,R^0f_*\mathbb Z).          \tag{4.3}
\]

\begin{theorem}[Engel's integral transgression]
Assume $|p|=1$, as in Engel's integral-homology calculation, and choose the
common orientation for which
$H^2(\mathbb P^1,R^0f_*\mathbb Z)=\mathbb Z\omega$.  Then
\[
 H^0(\mathbb P^1,R^1f_*\mathbb Z)=12\mathbb Z\psi,
 \qquad
                  d_2^{0,1}(12\psi)=p\omega.                      \tag{4.4}
\]
\end{theorem}

\begin{proof}
The extension lattice in (4.4) is Engel's Lemma~6.5.  For the differential,
the two local affine lifts satisfy
\[
                  \widetilde\rho_0^{\,3}=t_{v_0},
                  \qquad \widetilde\rho_1^{\,4}=t_{v_1}.
\]
Their oriented contributions to the obstruction for $12\psi$ are
\[
       \frac{12\psi(v_0)}3=4m,\qquad
       \frac{12\psi(v_1)}4=3n.                                   \tag{4.5}
\]
The cusp gluing at infinity contributes zero in this marking.  Their sum is
$p=4m+3n$, which is Engel's Lemma~6.6
\cite[pp.~17--18]{Engel}.  Reversing either generator in (4.4) reverses the
corresponding sign; simultaneous transport of the orientation preserves the
identity.
\end{proof}

\begin{proposition}[The $(3,4)$ presentation]
For the same geometric data,
\[
 \pi_1(Y(u))\cong
 \langle c,x_0,x_1\mid c\text{ central},\ x_0x_1=1,
          \ x_0^3=c^m,\ x_1^4=c^n\rangle.                         \tag{4.6}
\]
Consequently
\[
 \pi_1(Y(u))\cong
 \begin{cases}
  \mathbb Z/|p|\mathbb Z,&p\ne0,\\
  \mathbb Z,&p=0.
 \end{cases}                                                       \tag{4.7}
\]
In particular it is trivial when $|p|=1$.
\end{proposition}

\begin{proof}
Engel's vanishing cycles at infinity kill $\ker\psi$ and leave the central
class $c$ represented by $e_4$.  The two logarithmic fillings give (4.6)
\cite[Proposition 5.1, pp.~13--14]{Engel}.  This is Theorem~3.1 with
$(r,s;b_0)=(3,4;0)$, for which $d=1$ and $\Delta_0=p$.
\end{proof}

\subsection*{Limits of the coefficient}

The presentation theorem does not by itself define a map $f$, the sheaves
$R^bf_*\mathbb Z$, an extension lattice inside their generic invariants, or
a Leray differential.  It therefore gives no transgression theorem for
arbitrary $(r,s)$.

There is also a numerical obstruction to replacing $(3,4)$ formally by
arbitrary multiplicities while keeping the coefficient $rn+sm$.  In the
Seifert setting, let $g=\gcd(r,s)$ and $L=\operatorname{lcm}(r,s)$.  Equation
(3.13) gives
\[
 L e=-Lb_0+\frac{s}{g}m+\frac{r}{g}n
     =\frac{rn+sm-rsb_0}{g}.                                     \tag{4.8}
\]
Thus a transgression normalized by the denominator-clearing carrier $L\psi$
and by the rational Euler number would have coefficient
$\Delta_{b_0}/g$, not $\Delta_{b_0}$, unless $g=1$.  For the normalized
effective data
\[
 (r,s;\beta_0,\beta_1;b_0)=(2,2;1,1;0),
\]
one has $m=n=-1$, $L=2$, $e=-1$, and
\[
                   Le=-2\ne-4=\Delta_0.                           \tag{4.9}
\]
The global term $-Lb_0$ is another necessary contribution.  Engel's formula
avoids both ambiguities because $\gcd(3,4)=1$, $L=12$, the section gauge has
$b_0=0$, and the local extension lattice is computed geometrically.  A
general result requires equally explicit local models, specialization maps,
orientations, and a proof that the proposed carrier extends.

Finally, $\Lambda=H_1(F,\mathbb Z)$ is an auxiliary fibre lattice.  In the
case in which Engel's Theorem~3.1 identifies the total space with $S^6$, one
has
\[
 H^j(Y(u),\mathbb Z)=0\qquad(0<j<6),
\]
so neither $\Lambda$ nor its discriminant group is ordinary positive-degree
cohomology of the total space.

\begin{thebibliography}{9}
\raggedright

\bibitem{Engel}
P.~Engel,
\emph{Complex structures on $S^6$},
author's preprint, 2026,
\href{https://philip-engel.github.io/S6.pdf}{philip-engel.github.io/S6.pdf}.

\bibitem{Hatcher}
A.~Hatcher,
\emph{Notes on Basic $3$-Manifold Topology},
\href{https://pi.math.cornell.edu/~hatcher/3M/3M.pdf}
{pi.math.cornell.edu/$\sim$hatcher/3M/3M.pdf}.

\bibitem{Wilkes}
G.~Wilkes,
\emph{Profinite rigidity for Seifert fibre spaces},
Geom. Dedicata \textbf{188} (2017), 141--163,
\href{https://doi.org/10.1007/s10711-016-0209-6}
{doi:10.1007/s10711-016-0209-6}.

\end{thebibliography}

\end{document}
